\begin{table}[!htbp]\centering\scriptsize
\caption{Excess mass with minority interests frozen at end-2019 (panel 1 of 10)}
\begin{threeparttable}
\begin{tabular}{p{0.43\linewidth}lrrrr}
\toprule
Window, population, fit and ratio & Sample & N & Firms & b or difference & 95\% CI \\
\midrule
post T023 replication stored RL1 main official & All & 630 & 132 & 3.1429 & [1.0122, 5.8411] \\
post paired main official & All & 602 & 129 & 3.1818 & [1.0105, 6.0136] \\
post paired main frozen & All & 602 & 129 & -0.5250 & [-1.8293, 1.0763] \\
post paired main official minus frozen & All & 602 & 129 & 3.7068 & [1.6048, 6.1403] \\
post paired degree4 official & All & 602 & 129 & 2.9853 & [0.8507, 5.7669] \\
post paired degree4 frozen & All & 602 & 129 & -0.7028 & [-1.9565, 0.8282] \\
post paired degree4 official minus frozen & All & 602 & 129 & 3.6881 & [1.6542, 6.0124] \\
post paired degree7 official & All & 602 & 129 & 2.4876 & [0.5162, 5.1158] \\
post paired degree7 frozen & All & 602 & 129 & -0.6737 & [-1.8749, 0.9692] \\
post paired degree7 official minus frozen & All & 602 & 129 & 3.1612 & [1.1376, 5.5401] \\
\bottomrule
\end{tabular}
\begin{tablenotes}[flushleft]\footnotesize
\item[] These are mechanical accounting counterfactuals, not attainable balance sheets, causal effects, or economic leverage. Minority interests may be genuine equity. Official RL1 is N/(N+E), N=liabilities minus advances, E=total equity including minority interests. N and E must be positive. Frozen MI replaces current minority interests with their end-2019 level, holding all other observed items at their realized values; the pre-period placebo freezes them at end-2016. Ex-MI is mechanical, not a deconsolidated balance sheet. Nonpositive total counterfactual denominators are unavailable. All shares are proportions. N and firms count available observations and unique firms for each estimate; contributing contrast counts are identified in the table notes. Both variants are reported: all observations and exclude jump dates (known half-year minority-interest jumps above 50 percent of current equity). Unknown flags are retained; later dates retain the jump in the MI stock. Percentile 95 percent intervals use 999 whole-firm resamples with shared draws, full histories and complete re-estimation. Intervals are pointwise; no significance stars. HK changes assume unchanged reporting currency within firm; currency metadata remain unresolved. Interpretation rule specified before results: a large drop in b means bunching depends mechanically on post-2019 minority-interest growth; no drop means other items achieve bunching. The T023 stored-RL1 replication gate reproduces 3.14. The replication and all-eligible official fits share b despite different N because the 14 excluded nonpositive-N/E observations lie outside [0.40,1.00], leaving the same 557 histogram observations. Paired comparisons use identical eligible observations. All eligible comparisons retain all eligible official observations and available frozen observations, with separate counts and shared firm draws. Bins have width 0.01 over [0.40,1.00]; fifth-degree polynomial excludes [0.66,0.70). Alternatives use degree four or seven, [0.67,0.70), or [0.66,0.72). The mass window is always below 0.70. b equals excess count divided by the mean predicted count per mass-window bin. No integration constraint or clipping is used; nonpositive predicted counts invalidate a draw. Differences use paired draws with both fits re-estimated. Invalid draws and coverage are in the CSV; fewer than 90 percent valid draws makes inference unreliable.
\end{tablenotes}
\end{threeparttable}
\end{table}

\begin{table}[!htbp]\centering\scriptsize
\caption{Excess mass with minority interests frozen at end-2019 (panel 2 of 10)}
\begin{threeparttable}
\begin{tabular}{p{0.43\linewidth}lrrrr}
\toprule
Window, population, fit and ratio & Sample & N & Firms & b or difference & 95\% CI \\
\midrule
post paired narrow official & All & 602 & 129 & 3.0600 & [1.1420, 5.5682] \\
post paired narrow frozen & All & 602 & 129 & -0.4110 & [-1.4385, 0.9267] \\
post paired narrow official minus frozen & All & 602 & 129 & 3.4710 & [1.7327, 5.5173] \\
post paired wide official & All & 602 & 129 & 4.0215 & [1.4902, 7.1787] \\
post paired wide frozen & All & 602 & 129 & -0.2735 & [-1.6979, 1.6499] \\
post paired wide official minus frozen & All & 602 & 129 & 4.2950 & [1.9287, 7.0875] \\
post all eligible main official & All & 616 & 131 & 3.1429 & [1.0122, 5.8411] \\
post all eligible main frozen & All & 602 & 129 & -0.5250 & [-1.8293, 1.0763] \\
post all eligible main official minus frozen & All & 616 & 131 & 3.6678 & [1.5430, 5.9504] \\
post all eligible degree4 official & All & 616 & 131 & 2.9544 & [0.8841, 5.6043] \\
\bottomrule
\end{tabular}
\begin{tablenotes}[flushleft]\footnotesize
\item[] These are mechanical accounting counterfactuals, not attainable balance sheets, causal effects, or economic leverage. Minority interests may be genuine equity. Official RL1 is N/(N+E), N=liabilities minus advances, E=total equity including minority interests. N and E must be positive. Frozen MI replaces current minority interests with their end-2019 level, holding all other observed items at their realized values; the pre-period placebo freezes them at end-2016. Ex-MI is mechanical, not a deconsolidated balance sheet. Nonpositive total counterfactual denominators are unavailable. All shares are proportions. N and firms count available observations and unique firms for each estimate; contributing contrast counts are identified in the table notes. Both variants are reported: all observations and exclude jump dates (known half-year minority-interest jumps above 50 percent of current equity). Unknown flags are retained; later dates retain the jump in the MI stock. Percentile 95 percent intervals use 999 whole-firm resamples with shared draws, full histories and complete re-estimation. Intervals are pointwise; no significance stars. HK changes assume unchanged reporting currency within firm; currency metadata remain unresolved. Interpretation rule specified before results: a large drop in b means bunching depends mechanically on post-2019 minority-interest growth; no drop means other items achieve bunching. The T023 stored-RL1 replication gate reproduces 3.14. The replication and all-eligible official fits share b despite different N because the 14 excluded nonpositive-N/E observations lie outside [0.40,1.00], leaving the same 557 histogram observations. Paired comparisons use identical eligible observations. All eligible comparisons retain all eligible official observations and available frozen observations, with separate counts and shared firm draws. Bins have width 0.01 over [0.40,1.00]; fifth-degree polynomial excludes [0.66,0.70). Alternatives use degree four or seven, [0.67,0.70), or [0.66,0.72). The mass window is always below 0.70. b equals excess count divided by the mean predicted count per mass-window bin. No integration constraint or clipping is used; nonpositive predicted counts invalidate a draw. Differences use paired draws with both fits re-estimated. Invalid draws and coverage are in the CSV; fewer than 90 percent valid draws makes inference unreliable.
\end{tablenotes}
\end{threeparttable}
\end{table}

\begin{table}[!htbp]\centering\scriptsize
\caption{Excess mass with minority interests frozen at end-2019 (panel 3 of 10)}
\begin{threeparttable}
\begin{tabular}{p{0.43\linewidth}lrrrr}
\toprule
Window, population, fit and ratio & Sample & N & Firms & b or difference & 95\% CI \\
\midrule
post all eligible degree4 frozen & All & 602 & 129 & -0.7028 & [-1.9565, 0.8282] \\
post all eligible degree4 official minus frozen & All & 616 & 131 & 3.6572 & [1.5916, 5.8323] \\
post all eligible degree7 official & All & 616 & 131 & 2.4069 & [0.4372, 5.0671] \\
post all eligible degree7 frozen & All & 602 & 129 & -0.6737 & [-1.8749, 0.9692] \\
post all eligible degree7 official minus frozen & All & 616 & 131 & 3.0805 & [1.0673, 5.2910] \\
post all eligible narrow official & All & 616 & 131 & 2.9727 & [1.0747, 5.4119] \\
post all eligible narrow frozen & All & 602 & 129 & -0.4110 & [-1.4385, 0.9267] \\
post all eligible narrow official minus frozen & All & 616 & 131 & 3.3837 & [1.6702, 5.3167] \\
post all eligible wide official & All & 616 & 131 & 3.9735 & [1.4107, 7.1403] \\
post all eligible wide frozen & All & 602 & 129 & -0.2735 & [-1.6979, 1.6499] \\
\bottomrule
\end{tabular}
\begin{tablenotes}[flushleft]\footnotesize
\item[] These are mechanical accounting counterfactuals, not attainable balance sheets, causal effects, or economic leverage. Minority interests may be genuine equity. Official RL1 is N/(N+E), N=liabilities minus advances, E=total equity including minority interests. N and E must be positive. Frozen MI replaces current minority interests with their end-2019 level, holding all other observed items at their realized values; the pre-period placebo freezes them at end-2016. Ex-MI is mechanical, not a deconsolidated balance sheet. Nonpositive total counterfactual denominators are unavailable. All shares are proportions. N and firms count available observations and unique firms for each estimate; contributing contrast counts are identified in the table notes. Both variants are reported: all observations and exclude jump dates (known half-year minority-interest jumps above 50 percent of current equity). Unknown flags are retained; later dates retain the jump in the MI stock. Percentile 95 percent intervals use 999 whole-firm resamples with shared draws, full histories and complete re-estimation. Intervals are pointwise; no significance stars. HK changes assume unchanged reporting currency within firm; currency metadata remain unresolved. Interpretation rule specified before results: a large drop in b means bunching depends mechanically on post-2019 minority-interest growth; no drop means other items achieve bunching. The T023 stored-RL1 replication gate reproduces 3.14. The replication and all-eligible official fits share b despite different N because the 14 excluded nonpositive-N/E observations lie outside [0.40,1.00], leaving the same 557 histogram observations. Paired comparisons use identical eligible observations. All eligible comparisons retain all eligible official observations and available frozen observations, with separate counts and shared firm draws. Bins have width 0.01 over [0.40,1.00]; fifth-degree polynomial excludes [0.66,0.70). Alternatives use degree four or seven, [0.67,0.70), or [0.66,0.72). The mass window is always below 0.70. b equals excess count divided by the mean predicted count per mass-window bin. No integration constraint or clipping is used; nonpositive predicted counts invalidate a draw. Differences use paired draws with both fits re-estimated. Invalid draws and coverage are in the CSV; fewer than 90 percent valid draws makes inference unreliable.
\end{tablenotes}
\end{threeparttable}
\end{table}

\begin{table}[!htbp]\centering\scriptsize
\caption{Excess mass with minority interests frozen at end-2019 (panel 4 of 10)}
\begin{threeparttable}
\begin{tabular}{p{0.43\linewidth}lrrrr}
\toprule
Window, population, fit and ratio & Sample & N & Firms & b or difference & 95\% CI \\
\midrule
post all eligible wide official minus frozen & All & 616 & 131 & 4.2470 & [1.8484, 6.9876] \\
extension paired main official & All & 357 & 96 & 0.1405 & [-1.8019, 2.4956] \\
extension paired main frozen & All & 357 & 96 & 0.1613 & [-2.0230, 2.9137] \\
extension paired main official minus frozen & All & 357 & 96 & -0.0208 & [-2.4852, 2.2439] \\
extension paired degree4 official & All & 357 & 96 & 0.2475 & [-1.7300, 2.7284] \\
extension paired degree4 frozen & All & 357 & 96 & 0.1864 & [-2.0338, 3.0738] \\
extension paired degree4 official minus frozen & All & 357 & 96 & 0.0611 & [-2.4704, 2.3295] \\
extension paired degree7 official & All & 357 & 96 & -0.2446 & [-2.0029, 1.9102] \\
extension paired degree7 frozen & All & 357 & 96 & 0.0436 & [-2.0653, 3.3285] \\
extension paired degree7 official minus frozen & All & 357 & 96 & -0.2882 & [-2.9010, 1.9157] \\
\bottomrule
\end{tabular}
\begin{tablenotes}[flushleft]\footnotesize
\item[] These are mechanical accounting counterfactuals, not attainable balance sheets, causal effects, or economic leverage. Minority interests may be genuine equity. Official RL1 is N/(N+E), N=liabilities minus advances, E=total equity including minority interests. N and E must be positive. Frozen MI replaces current minority interests with their end-2019 level, holding all other observed items at their realized values; the pre-period placebo freezes them at end-2016. Ex-MI is mechanical, not a deconsolidated balance sheet. Nonpositive total counterfactual denominators are unavailable. All shares are proportions. N and firms count available observations and unique firms for each estimate; contributing contrast counts are identified in the table notes. Both variants are reported: all observations and exclude jump dates (known half-year minority-interest jumps above 50 percent of current equity). Unknown flags are retained; later dates retain the jump in the MI stock. Percentile 95 percent intervals use 999 whole-firm resamples with shared draws, full histories and complete re-estimation. Intervals are pointwise; no significance stars. HK changes assume unchanged reporting currency within firm; currency metadata remain unresolved. Interpretation rule specified before results: a large drop in b means bunching depends mechanically on post-2019 minority-interest growth; no drop means other items achieve bunching. The T023 stored-RL1 replication gate reproduces 3.14. The replication and all-eligible official fits share b despite different N because the 14 excluded nonpositive-N/E observations lie outside [0.40,1.00], leaving the same 557 histogram observations. Paired comparisons use identical eligible observations. All eligible comparisons retain all eligible official observations and available frozen observations, with separate counts and shared firm draws. Bins have width 0.01 over [0.40,1.00]; fifth-degree polynomial excludes [0.66,0.70). Alternatives use degree four or seven, [0.67,0.70), or [0.66,0.72). The mass window is always below 0.70. b equals excess count divided by the mean predicted count per mass-window bin. No integration constraint or clipping is used; nonpositive predicted counts invalidate a draw. Differences use paired draws with both fits re-estimated. Invalid draws and coverage are in the CSV; fewer than 90 percent valid draws makes inference unreliable.
\end{tablenotes}
\end{threeparttable}
\end{table}

\begin{table}[!htbp]\centering\scriptsize
\caption{Excess mass with minority interests frozen at end-2019 (panel 5 of 10)}
\begin{threeparttable}
\begin{tabular}{p{0.43\linewidth}lrrrr}
\toprule
Window, population, fit and ratio & Sample & N & Firms & b or difference & 95\% CI \\
\midrule
extension paired narrow official & All & 357 & 96 & -0.2816 & [-1.6561, 1.3395] \\
extension paired narrow frozen & All & 357 & 96 & 0.3371 & [-1.4569, 2.5356] \\
extension paired narrow official minus frozen & All & 357 & 96 & -0.6187 & [-2.4724, 1.0694] \\
extension paired wide official & All & 357 & 96 & 0.2999 & [-1.7905, 2.9737] \\
extension paired wide frozen & All & 357 & 96 & 0.1031 & [-2.0580, 3.0948] \\
extension paired wide official minus frozen & All & 357 & 96 & 0.1968 & [-2.1346, 2.4624] \\
post paired main official & Excl. dates & 592 & 129 & 3.1625 & [0.9975, 6.0070] \\
post paired main frozen & Excl. dates & 592 & 129 & -0.5281 & [-1.8344, 1.0552] \\
post paired main official minus frozen & Excl. dates & 592 & 129 & 3.6906 & [1.5945, 6.0875] \\
post paired degree4 official & Excl. dates & 592 & 129 & 2.9610 & [0.8403, 5.7148] \\
\bottomrule
\end{tabular}
\begin{tablenotes}[flushleft]\footnotesize
\item[] These are mechanical accounting counterfactuals, not attainable balance sheets, causal effects, or economic leverage. Minority interests may be genuine equity. Official RL1 is N/(N+E), N=liabilities minus advances, E=total equity including minority interests. N and E must be positive. Frozen MI replaces current minority interests with their end-2019 level, holding all other observed items at their realized values; the pre-period placebo freezes them at end-2016. Ex-MI is mechanical, not a deconsolidated balance sheet. Nonpositive total counterfactual denominators are unavailable. All shares are proportions. N and firms count available observations and unique firms for each estimate; contributing contrast counts are identified in the table notes. Both variants are reported: all observations and exclude jump dates (known half-year minority-interest jumps above 50 percent of current equity). Unknown flags are retained; later dates retain the jump in the MI stock. Percentile 95 percent intervals use 999 whole-firm resamples with shared draws, full histories and complete re-estimation. Intervals are pointwise; no significance stars. HK changes assume unchanged reporting currency within firm; currency metadata remain unresolved. Interpretation rule specified before results: a large drop in b means bunching depends mechanically on post-2019 minority-interest growth; no drop means other items achieve bunching. The T023 stored-RL1 replication gate reproduces 3.14. The replication and all-eligible official fits share b despite different N because the 14 excluded nonpositive-N/E observations lie outside [0.40,1.00], leaving the same 557 histogram observations. Paired comparisons use identical eligible observations. All eligible comparisons retain all eligible official observations and available frozen observations, with separate counts and shared firm draws. Bins have width 0.01 over [0.40,1.00]; fifth-degree polynomial excludes [0.66,0.70). Alternatives use degree four or seven, [0.67,0.70), or [0.66,0.72). The mass window is always below 0.70. b equals excess count divided by the mean predicted count per mass-window bin. No integration constraint or clipping is used; nonpositive predicted counts invalidate a draw. Differences use paired draws with both fits re-estimated. Invalid draws and coverage are in the CSV; fewer than 90 percent valid draws makes inference unreliable.
\end{tablenotes}
\end{threeparttable}
\end{table}

\begin{table}[!htbp]\centering\scriptsize
\caption{Excess mass with minority interests frozen at end-2019 (panel 6 of 10)}
\begin{threeparttable}
\begin{tabular}{p{0.43\linewidth}lrrrr}
\toprule
Window, population, fit and ratio & Sample & N & Firms & b or difference & 95\% CI \\
\midrule
post paired degree4 frozen & Excl. dates & 592 & 129 & -0.7069 & [-1.9571, 0.8310] \\
post paired degree4 official minus frozen & Excl. dates & 592 & 129 & 3.6679 & [1.6327, 5.9610] \\
post paired degree7 official & Excl. dates & 592 & 129 & 2.4723 & [0.5103, 5.0880] \\
post paired degree7 frozen & Excl. dates & 592 & 129 & -0.6865 & [-1.8850, 0.9362] \\
post paired degree7 official minus frozen & Excl. dates & 592 & 129 & 3.1588 & [1.1477, 5.5159] \\
post paired narrow official & Excl. dates & 592 & 129 & 3.0452 & [1.1253, 5.5226] \\
post paired narrow frozen & Excl. dates & 592 & 129 & -0.4119 & [-1.4371, 0.9150] \\
post paired narrow official minus frozen & Excl. dates & 592 & 129 & 3.4571 & [1.7296, 5.5026] \\
post paired wide official & Excl. dates & 592 & 129 & 3.9942 & [1.4746, 7.1216] \\
post paired wide frozen & Excl. dates & 592 & 129 & -0.2755 & [-1.6976, 1.6531] \\
\bottomrule
\end{tabular}
\begin{tablenotes}[flushleft]\footnotesize
\item[] These are mechanical accounting counterfactuals, not attainable balance sheets, causal effects, or economic leverage. Minority interests may be genuine equity. Official RL1 is N/(N+E), N=liabilities minus advances, E=total equity including minority interests. N and E must be positive. Frozen MI replaces current minority interests with their end-2019 level, holding all other observed items at their realized values; the pre-period placebo freezes them at end-2016. Ex-MI is mechanical, not a deconsolidated balance sheet. Nonpositive total counterfactual denominators are unavailable. All shares are proportions. N and firms count available observations and unique firms for each estimate; contributing contrast counts are identified in the table notes. Both variants are reported: all observations and exclude jump dates (known half-year minority-interest jumps above 50 percent of current equity). Unknown flags are retained; later dates retain the jump in the MI stock. Percentile 95 percent intervals use 999 whole-firm resamples with shared draws, full histories and complete re-estimation. Intervals are pointwise; no significance stars. HK changes assume unchanged reporting currency within firm; currency metadata remain unresolved. Interpretation rule specified before results: a large drop in b means bunching depends mechanically on post-2019 minority-interest growth; no drop means other items achieve bunching. The T023 stored-RL1 replication gate reproduces 3.14. The replication and all-eligible official fits share b despite different N because the 14 excluded nonpositive-N/E observations lie outside [0.40,1.00], leaving the same 557 histogram observations. Paired comparisons use identical eligible observations. All eligible comparisons retain all eligible official observations and available frozen observations, with separate counts and shared firm draws. Bins have width 0.01 over [0.40,1.00]; fifth-degree polynomial excludes [0.66,0.70). Alternatives use degree four or seven, [0.67,0.70), or [0.66,0.72). The mass window is always below 0.70. b equals excess count divided by the mean predicted count per mass-window bin. No integration constraint or clipping is used; nonpositive predicted counts invalidate a draw. Differences use paired draws with both fits re-estimated. Invalid draws and coverage are in the CSV; fewer than 90 percent valid draws makes inference unreliable.
\end{tablenotes}
\end{threeparttable}
\end{table}

\begin{table}[!htbp]\centering\scriptsize
\caption{Excess mass with minority interests frozen at end-2019 (panel 7 of 10)}
\begin{threeparttable}
\begin{tabular}{p{0.43\linewidth}lrrrr}
\toprule
Window, population, fit and ratio & Sample & N & Firms & b or difference & 95\% CI \\
\midrule
post paired wide official minus frozen & Excl. dates & 592 & 129 & 4.2697 & [1.9146, 7.0326] \\
post all eligible main official & Excl. dates & 606 & 131 & 3.1241 & [1.0003, 5.8166] \\
post all eligible main frozen & Excl. dates & 592 & 129 & -0.5281 & [-1.8344, 1.0552] \\
post all eligible main official minus frozen & Excl. dates & 606 & 131 & 3.6522 & [1.5418, 5.9159] \\
post all eligible degree4 official & Excl. dates & 606 & 131 & 2.9308 & [0.8686, 5.6044] \\
post all eligible degree4 frozen & Excl. dates & 592 & 129 & -0.7069 & [-1.9571, 0.8310] \\
post all eligible degree4 official minus frozen & Excl. dates & 606 & 131 & 3.6376 & [1.5858, 5.7856] \\
post all eligible degree7 official & Excl. dates & 606 & 131 & 2.3922 & [0.4283, 5.0428] \\
post all eligible degree7 frozen & Excl. dates & 592 & 129 & -0.6865 & [-1.8850, 0.9362] \\
post all eligible degree7 official minus frozen & Excl. dates & 606 & 131 & 3.0787 & [1.0839, 5.2838] \\
\bottomrule
\end{tabular}
\begin{tablenotes}[flushleft]\footnotesize
\item[] These are mechanical accounting counterfactuals, not attainable balance sheets, causal effects, or economic leverage. Minority interests may be genuine equity. Official RL1 is N/(N+E), N=liabilities minus advances, E=total equity including minority interests. N and E must be positive. Frozen MI replaces current minority interests with their end-2019 level, holding all other observed items at their realized values; the pre-period placebo freezes them at end-2016. Ex-MI is mechanical, not a deconsolidated balance sheet. Nonpositive total counterfactual denominators are unavailable. All shares are proportions. N and firms count available observations and unique firms for each estimate; contributing contrast counts are identified in the table notes. Both variants are reported: all observations and exclude jump dates (known half-year minority-interest jumps above 50 percent of current equity). Unknown flags are retained; later dates retain the jump in the MI stock. Percentile 95 percent intervals use 999 whole-firm resamples with shared draws, full histories and complete re-estimation. Intervals are pointwise; no significance stars. HK changes assume unchanged reporting currency within firm; currency metadata remain unresolved. Interpretation rule specified before results: a large drop in b means bunching depends mechanically on post-2019 minority-interest growth; no drop means other items achieve bunching. The T023 stored-RL1 replication gate reproduces 3.14. The replication and all-eligible official fits share b despite different N because the 14 excluded nonpositive-N/E observations lie outside [0.40,1.00], leaving the same 557 histogram observations. Paired comparisons use identical eligible observations. All eligible comparisons retain all eligible official observations and available frozen observations, with separate counts and shared firm draws. Bins have width 0.01 over [0.40,1.00]; fifth-degree polynomial excludes [0.66,0.70). Alternatives use degree four or seven, [0.67,0.70), or [0.66,0.72). The mass window is always below 0.70. b equals excess count divided by the mean predicted count per mass-window bin. No integration constraint or clipping is used; nonpositive predicted counts invalidate a draw. Differences use paired draws with both fits re-estimated. Invalid draws and coverage are in the CSV; fewer than 90 percent valid draws makes inference unreliable.
\end{tablenotes}
\end{threeparttable}
\end{table}

\begin{table}[!htbp]\centering\scriptsize
\caption{Excess mass with minority interests frozen at end-2019 (panel 8 of 10)}
\begin{threeparttable}
\begin{tabular}{p{0.43\linewidth}lrrrr}
\toprule
Window, population, fit and ratio & Sample & N & Firms & b or difference & 95\% CI \\
\midrule
post all eligible narrow official & Excl. dates & 606 & 131 & 2.9584 & [1.0628, 5.3851] \\
post all eligible narrow frozen & Excl. dates & 592 & 129 & -0.4119 & [-1.4371, 0.9150] \\
post all eligible narrow official minus frozen & Excl. dates & 606 & 131 & 3.3703 & [1.6650, 5.2938] \\
post all eligible wide official & Excl. dates & 606 & 131 & 3.9470 & [1.3819, 7.0885] \\
post all eligible wide frozen & Excl. dates & 592 & 129 & -0.2755 & [-1.6976, 1.6531] \\
post all eligible wide official minus frozen & Excl. dates & 606 & 131 & 4.2225 & [1.8339, 6.9482] \\
extension paired main official & Excl. dates & 354 & 96 & 0.1376 & [-1.8054, 2.4908] \\
extension paired main frozen & Excl. dates & 354 & 96 & 0.1483 & [-2.0319, 2.8989] \\
extension paired main official minus frozen & Excl. dates & 354 & 96 & -0.0106 & [-2.4479, 2.2406] \\
extension paired degree4 official & Excl. dates & 354 & 96 & 0.2448 & [-1.7424, 2.7315] \\
\bottomrule
\end{tabular}
\begin{tablenotes}[flushleft]\footnotesize
\item[] These are mechanical accounting counterfactuals, not attainable balance sheets, causal effects, or economic leverage. Minority interests may be genuine equity. Official RL1 is N/(N+E), N=liabilities minus advances, E=total equity including minority interests. N and E must be positive. Frozen MI replaces current minority interests with their end-2019 level, holding all other observed items at their realized values; the pre-period placebo freezes them at end-2016. Ex-MI is mechanical, not a deconsolidated balance sheet. Nonpositive total counterfactual denominators are unavailable. All shares are proportions. N and firms count available observations and unique firms for each estimate; contributing contrast counts are identified in the table notes. Both variants are reported: all observations and exclude jump dates (known half-year minority-interest jumps above 50 percent of current equity). Unknown flags are retained; later dates retain the jump in the MI stock. Percentile 95 percent intervals use 999 whole-firm resamples with shared draws, full histories and complete re-estimation. Intervals are pointwise; no significance stars. HK changes assume unchanged reporting currency within firm; currency metadata remain unresolved. Interpretation rule specified before results: a large drop in b means bunching depends mechanically on post-2019 minority-interest growth; no drop means other items achieve bunching. The T023 stored-RL1 replication gate reproduces 3.14. The replication and all-eligible official fits share b despite different N because the 14 excluded nonpositive-N/E observations lie outside [0.40,1.00], leaving the same 557 histogram observations. Paired comparisons use identical eligible observations. All eligible comparisons retain all eligible official observations and available frozen observations, with separate counts and shared firm draws. Bins have width 0.01 over [0.40,1.00]; fifth-degree polynomial excludes [0.66,0.70). Alternatives use degree four or seven, [0.67,0.70), or [0.66,0.72). The mass window is always below 0.70. b equals excess count divided by the mean predicted count per mass-window bin. No integration constraint or clipping is used; nonpositive predicted counts invalidate a draw. Differences use paired draws with both fits re-estimated. Invalid draws and coverage are in the CSV; fewer than 90 percent valid draws makes inference unreliable.
\end{tablenotes}
\end{threeparttable}
\end{table}

\begin{table}[!htbp]\centering\scriptsize
\caption{Excess mass with minority interests frozen at end-2019 (panel 9 of 10)}
\begin{threeparttable}
\begin{tabular}{p{0.43\linewidth}lrrrr}
\toprule
Window, population, fit and ratio & Sample & N & Firms & b or difference & 95\% CI \\
\midrule
extension paired degree4 frozen & Excl. dates & 354 & 96 & 0.1626 & [-2.0419, 3.0361] \\
extension paired degree4 official minus frozen & Excl. dates & 354 & 96 & 0.0821 & [-2.4038, 2.3452] \\
extension paired degree7 official & Excl. dates & 354 & 96 & -0.2430 & [-2.0023, 1.9063] \\
extension paired degree7 frozen & Excl. dates & 354 & 96 & 0.0444 & [-2.0626, 3.3564] \\
extension paired degree7 official minus frozen & Excl. dates & 354 & 96 & -0.2874 & [-2.9100, 1.9129] \\
extension paired narrow official & Excl. dates & 354 & 96 & -0.2834 & [-1.6599, 1.3475] \\
extension paired narrow frozen & Excl. dates & 354 & 96 & 0.3262 & [-1.4571, 2.5093] \\
extension paired narrow official minus frozen & Excl. dates & 354 & 96 & -0.6096 & [-2.4305, 1.0795] \\
extension paired wide official & Excl. dates & 354 & 96 & 0.2961 & [-1.7964, 3.0029] \\
extension paired wide frozen & Excl. dates & 354 & 96 & 0.0874 & [-2.0654, 3.0763] \\
\bottomrule
\end{tabular}
\begin{tablenotes}[flushleft]\footnotesize
\item[] These are mechanical accounting counterfactuals, not attainable balance sheets, causal effects, or economic leverage. Minority interests may be genuine equity. Official RL1 is N/(N+E), N=liabilities minus advances, E=total equity including minority interests. N and E must be positive. Frozen MI replaces current minority interests with their end-2019 level, holding all other observed items at their realized values; the pre-period placebo freezes them at end-2016. Ex-MI is mechanical, not a deconsolidated balance sheet. Nonpositive total counterfactual denominators are unavailable. All shares are proportions. N and firms count available observations and unique firms for each estimate; contributing contrast counts are identified in the table notes. Both variants are reported: all observations and exclude jump dates (known half-year minority-interest jumps above 50 percent of current equity). Unknown flags are retained; later dates retain the jump in the MI stock. Percentile 95 percent intervals use 999 whole-firm resamples with shared draws, full histories and complete re-estimation. Intervals are pointwise; no significance stars. HK changes assume unchanged reporting currency within firm; currency metadata remain unresolved. Interpretation rule specified before results: a large drop in b means bunching depends mechanically on post-2019 minority-interest growth; no drop means other items achieve bunching. The T023 stored-RL1 replication gate reproduces 3.14. The replication and all-eligible official fits share b despite different N because the 14 excluded nonpositive-N/E observations lie outside [0.40,1.00], leaving the same 557 histogram observations. Paired comparisons use identical eligible observations. All eligible comparisons retain all eligible official observations and available frozen observations, with separate counts and shared firm draws. Bins have width 0.01 over [0.40,1.00]; fifth-degree polynomial excludes [0.66,0.70). Alternatives use degree four or seven, [0.67,0.70), or [0.66,0.72). The mass window is always below 0.70. b equals excess count divided by the mean predicted count per mass-window bin. No integration constraint or clipping is used; nonpositive predicted counts invalidate a draw. Differences use paired draws with both fits re-estimated. Invalid draws and coverage are in the CSV; fewer than 90 percent valid draws makes inference unreliable.
\end{tablenotes}
\end{threeparttable}
\end{table}

\begin{table}[!htbp]\centering\scriptsize
\caption{Excess mass with minority interests frozen at end-2019 (panel 10 of 10)}
\begin{threeparttable}
\begin{tabular}{p{0.43\linewidth}lrrrr}
\toprule
Window, population, fit and ratio & Sample & N & Firms & b or difference & 95\% CI \\
\midrule
extension paired wide official minus frozen & Excl. dates & 354 & 96 & 0.2087 & [-2.1260, 2.4744] \\
\bottomrule
\end{tabular}
\begin{tablenotes}[flushleft]\footnotesize
\item[] These are mechanical accounting counterfactuals, not attainable balance sheets, causal effects, or economic leverage. Minority interests may be genuine equity. Official RL1 is N/(N+E), N=liabilities minus advances, E=total equity including minority interests. N and E must be positive. Frozen MI replaces current minority interests with their end-2019 level, holding all other observed items at their realized values; the pre-period placebo freezes them at end-2016. Ex-MI is mechanical, not a deconsolidated balance sheet. Nonpositive total counterfactual denominators are unavailable. All shares are proportions. N and firms count available observations and unique firms for each estimate; contributing contrast counts are identified in the table notes. Both variants are reported: all observations and exclude jump dates (known half-year minority-interest jumps above 50 percent of current equity). Unknown flags are retained; later dates retain the jump in the MI stock. Percentile 95 percent intervals use 999 whole-firm resamples with shared draws, full histories and complete re-estimation. Intervals are pointwise; no significance stars. HK changes assume unchanged reporting currency within firm; currency metadata remain unresolved. Interpretation rule specified before results: a large drop in b means bunching depends mechanically on post-2019 minority-interest growth; no drop means other items achieve bunching. The T023 stored-RL1 replication gate reproduces 3.14. The replication and all-eligible official fits share b despite different N because the 14 excluded nonpositive-N/E observations lie outside [0.40,1.00], leaving the same 557 histogram observations. Paired comparisons use identical eligible observations. All eligible comparisons retain all eligible official observations and available frozen observations, with separate counts and shared firm draws. Bins have width 0.01 over [0.40,1.00]; fifth-degree polynomial excludes [0.66,0.70). Alternatives use degree four or seven, [0.67,0.70), or [0.66,0.72). The mass window is always below 0.70. b equals excess count divided by the mean predicted count per mass-window bin. No integration constraint or clipping is used; nonpositive predicted counts invalidate a draw. Differences use paired draws with both fits re-estimated. Invalid draws and coverage are in the CSV; fewer than 90 percent valid draws makes inference unreliable.
\end{tablenotes}
\end{threeparttable}
\end{table}

